\SourcePage{463}{466}
\section{Selection rules for matrix elements}

Group theory serves not only to classify the terms of any symmetric physical
system; it also supplies a straightforward way of determining selection rules
for matrix elements of the various quantities that characterize the system.

The method rests on the following general theorem. Let
$\psi_i^{(\alpha)}$ be one of the functions in a basis of an irreducible
(nonidentity) representation of the symmetry group. Its integral over the
entire space\footnote{The configuration space of the physical system in
question is meant here.} then vanishes identically:
\begin{equation}
 \int \psi_i^{(\alpha)}\,dq=0.
 \tag{97.1}
\end{equation}
To prove this, observe that an integral over all space is invariant under any
change of the coordinate system and, in particular,
\SourcePage{464}{467}
under every symmetry transformation. Hence
\[
 \int\psi_i^{(\alpha)}dq=\int\widehat G\psi_i^{(\alpha)}dq
 =\int\sum_k G_{ki}^{(\alpha)}\psi_k^{(\alpha)}dq.
\]
Now sum this equality over the elements of the group. On the left the integral
is merely multiplied by the group order $g$, so that
\[
 g\int\psi_i^{(\alpha)}dq=\sum_k\int\psi_k^{(\alpha)}\sum_G G_{ki}^{(\alpha)}dq.
\]
For any irreducible representation other than the identity representation,
however, $\sum_GG_{ki}^{(\alpha)}=0$ identically. This is the special case of
the orthogonality relations (94.7) in which one of the two irreducible
representations is the identity representation. The theorem follows.

If $\psi$ is a basis function of some reducible representation, the integral
$\int\psi\,dq$ can be nonzero only when that representation includes the
identity representation. This assertion is an immediate consequence of the
preceding theorem.

The matrix elements of a physical quantity $f$ are the integrals
\begin{equation}
 \langle\beta k|f|\alpha i\rangle=\int\psi_k^{(\beta)}\widehat f\psi_i^{(\alpha)}dq,
 \tag{97.2}
\end{equation}
where $\alpha$ and $\beta$ distinguish different energy levels of the system,
while $i$ and $k$ enumerate the wave functions belonging to the same
degenerate level\footnote{After passing to ``physically irreducible''
representations, the basis functions may be chosen real. For this reason,
(97.2) makes no distinction between the wave functions and their complex
conjugates.}. Denote by $D^{(\alpha)}$ and $D^{(\beta)}$ the irreducible
representations of the system's symmetry group realized by
$\psi_i^{(\alpha)}$ and $\psi_k^{(\beta)}$. Let $D_f$ denote the
representation of that group corresponding to the symmetry of $f$; it is
determined by the tensor character of $f$. Thus, when $f$ is a true scalar,
its operator $\widehat f$ is invariant under every symmetry transformation,
and $D_f$ is the identity representation. The same conclusion holds for a
pseudoscalar if the group contains symmetry axes only. If reflections also
belong to the group, $D_f$ is one-dimensional but is not the identity
representation. For a vector quantity $f$, the representation $D_f$ is
realized by the three vector components that transform into one another;
\SourcePage{465}{468}
in general this representation is different for polar and axial vectors.

The products $\psi_k^{(\beta)}\widehat f\psi_i^{(\alpha)}$ realize the group
representation given by the direct product
$D^{(\beta)}\times D_f\times D^{(\alpha)}$. A matrix element can be nonzero
when this representation contains the identity representation, or,
equivalently, when $D^{(\beta)}\times D^{(\alpha)}$ contains $D_f$. In
practice it is more convenient to resolve $D^{(\alpha)}\times D_f$ into
irreducible parts. This immediately gives every type $D^{(\beta)}$ of state
for which transitions from a state of type $D^{(\alpha)}$ can have nonzero
matrix elements.

The simplest case is that of a scalar quantity, for which $D_f$ is the
identity representation. It follows at once that a matrix element can be
nonzero only for a transition between states of the same type. Indeed, the
direct product of two distinct irreducible representations
$D^{(\alpha)}\times D^{(\beta)}$ does not contain the identity
representation, whereas the direct product of an irreducible representation
with itself always does. This is the most general statement of a theorem whose
special cases have already arisen several times.

Matrix elements diagonal in energy call for separate treatment: they describe transitions between states belonging to one and the same term, in contrast to transitions between states belonging to two distinct terms of the same type.

Here there is only one system of functions
$\psi_1^{(\alpha)}$, $\psi_2^{(\alpha)}$,~\ldots, rather than two different
systems. The procedure for obtaining the selection rules now depends on how
$f$ behaves under time reversal.

Consider a state whose wave function has the form
$\psi=\sum_i c_i\psi_i^{(\alpha)}$. In this state the mean value of $f$ is
\[
 \bar f=\sum_{i,k}c_k^*c_i\langle\alpha k|f|\alpha i\rangle.
\]
For the state with the complex-conjugate wave function
$\psi^*=\sum_i c_i^*\psi_i^{(\alpha)}$, we have
\[
 \bar f=\sum_{i,k}c_kc_i^*\langle\alpha k|f|\alpha i\rangle
 =\sum_{i,k}c_ic_k^*\langle\alpha i|f|\alpha k\rangle.
\]
When $f$ is unchanged by time reversal, the two states share not only their energy level

\SourcePage{466}{469}
but also the same value $\bar f$.
Since the coefficients $c_i$ are arbitrary, one must therefore have:

\[
 \langle\alpha k|f|\alpha i\rangle=\langle\alpha i|f|\alpha k\rangle.
\]
It is readily shown that the selection rules must then be found from the
symmetric part $[D^{(\alpha)2}]$ rather than from the whole direct product
$D^{(\alpha)}\times D^{(\alpha)}$. Nonzero matrix elements exist if
$[D^{(\alpha)2}]$ contains $D_f$\footnote{The product
$[D^{(\alpha)2}]$ always includes the identity representation. Thus, for a
scalar quantity, diagonal elements (as well as nondiagonal elements between
states of the same type) can be nonzero.}.

If instead $f$ changes sign under time reversal, replacing $\psi$ by $\psi^*$
must be accompanied by a reversal of the sign of $\bar f$. The same argument
then yields
\[
 \langle\alpha k|f|\alpha i\rangle=-\langle\alpha i|f|\alpha k\rangle.
\]
In this case the selection rules are fixed by the decomposition of the
antisymmetric part of the direct product, $\{D^{(\alpha)2}\}$.

\subsection*{Problems}

\ProblemHeading{1.} Find the selection rules for the matrix elements of the
electric dipole moment $\mathbf d$ and magnetic dipole moment
$\boldsymbol\mu$ when the symmetry is $O$.

\SolutionHeading The group $O$ contains no reflections. Polar vectors
($\mathbf d$) and axial vectors ($\boldsymbol\mu$) therefore transform
according to the same irreducible representation, $F_1$. The direct products
of $F_1$ with the other representations of $O$ decompose as
\begin{align}
 F_1\times A_1&=F_2,&F_1\times A_2&=F_1,&F_1\times E&=F_1+F_2,\notag\\
 F_1\times F_1&=A_1+E+F_1+F_2,&F_1\times F_2&=A_2+E+F_1+F_2.\tag{1}
\end{align}
Consequently, the nondiagonal (in energy) matrix elements can be nonzero for
the transitions
\[
 F_1\leftrightarrow A_1,E,F_1,F_2;\qquad F_2\leftrightarrow A_2,E,F_1,F_2.
\]
The symmetric and antisymmetric products of the irreducible representations
of $O$ are
\begin{align}
 [A_1^2]&=[A_2^2]=A_1,&[E^2]&=A_1+E,&[F_1^2]&=[F_2^2]=A_1+E+F_2,\notag\\
 \{E^2\}&=A_2,&\{F_1^2\}&=\{F_2^2\}=F_1.\tag{2}
\end{align}
None of the symmetric products contains $F_1$. Hence the vector $\mathbf d$,
which is invariant under time reversal, has no matrix elements diagonal in
energy. The magnetic moment, on the other hand, reverses sign under time
reversal and has diagonal matrix elements in states of type $F_1$ and $F_2$.

\ProblemHeading{2.} Repeat the calculation for symmetry $D_{3d}$.

\SolutionHeading The transformation laws of $\mathbf d$ and
$\boldsymbol\mu$ differ in the group $D_{3d}$:
\[
 d_x,d_y\sim E_u,\quad d_z\sim A_{2u},\qquad
 \mu_x,\mu_y\sim E_g,\quad \mu_z\sim A_{2g}
\]
\SourcePage{467}{470}
(here and in the problems below, the sign $\sim$ stands for the words
``transforms according to the representation''). We have
\begin{align}
 E_u\times A_{1g}&=E_u\times A_{2g}=E_u,&E_u\times A_{1u}&=E_u\times A_{2u}=E_g,\notag\\
 E_u\times E_u&=A_{1g}+A_{2g}+E_g,&E_u\times E_g&=A_{1u}+A_{2u}+E_u.\tag{3}
\end{align}
Thus the nondiagonal matrix elements of $d_x,d_y$ can be nonzero for
$E_u\leftrightarrow A_{1g},A_{2g},E_g$ and
$E_g\leftrightarrow A_{1u},A_{2u},E_u$. The other selection rules are found
in the same way:
\begin{align*}
\text{for }d_z:\quad&A_{1g}\leftrightarrow A_{2u};\ A_{2g}\leftrightarrow A_{1u};\ E_g\leftrightarrow E_u;\\
\text{for }\mu_x,\mu_y:\quad&E_g\leftrightarrow A_{1g},A_{2g},E_g;\ E_u\leftrightarrow A_{1u},A_{2u},E_u;\\
\text{for }\mu_z:\quad&A_{1g}\leftrightarrow A_{2g};\ A_{1u}\leftrightarrow A_{2u};\ E_g\leftrightarrow E_g;\ E_u\leftrightarrow E_u.
\end{align*}
The symmetric and antisymmetric products of the irreducible representations
of $D_{3d}$ are
\begin{align}
 [A_{1g}^2]&=[A_{1u}^2]=[A_{2g}^2]=[A_{2u}^2]=A_{1g},\notag\\
 [E_g^2]&=[E_u^2]=E_g+A_{1g},\notag\\
 \{E_g^2\}&=\{E_u^2\}=A_{2g}.\tag{4}
\end{align}
It follows that every component of $\mathbf d$ has vanishing matrix elements
diagonal in energy. For $\boldsymbol\mu$, diagonal matrix elements of $\mu_z$
exist for transitions between states belonging to a degenerate level of type
$E_g$ or $E_u$.

\ProblemHeading{3.} Find the selection rules for matrix elements of the
electric quadrupole-moment tensor $Q_{ik}$ for symmetry $O$.

\SolutionHeading With respect to the group $O$, the components of $Q_{ik}$,
a symmetric tensor satisfying $\sum_iQ_{ii}=0$, transform as follows:
\[
 Q_{xy},\ Q_{xz},\ Q_{yz}\sim F_2,
\]
\[
 Q_{xx}+\epsilon Q_{yy}+\epsilon^2Q_{zz},\quad Q_{xx}+\epsilon^2Q_{yy}+\epsilon Q_{zz}\sim E
 \quad(\epsilon=e^{2\pi i/3}).
\]
Decomposing the direct products of $F_2$ and $E$ with every representation of
the group gives the following selection rules for nondiagonal matrix elements:
\begin{align*}
 \text{for }Q_{xy},Q_{xz},Q_{yz}:&\quad F_1\leftrightarrow A_2,E,F_1,F_2;\ F_2\leftrightarrow A_1,E,F_1,F_2;\\
 \text{for }Q_{xx},Q_{yy},Q_{zz}:&\quad E\leftrightarrow A_1,A_2,E;\ F_1\leftrightarrow F_1,F_2;\ F_2\leftrightarrow F_2.
\end{align*}
As follows from (2), diagonal matrix elements occur in the following states:
\begin{align*}
 \text{for }Q_{xy},Q_{xz},Q_{yz}:&\quad F_1,F_2,\\
 \text{for }Q_{xx},Q_{yy},Q_{zz}:&\quad E,F_1,F_2.
\end{align*}

\ProblemHeading{4.} Repeat the calculation for symmetry $D_{3d}$.

\SolutionHeading The components of $Q_{ik}$ transform under $D_{3d}$
according to
\[
 Q_{zz}\sim A_{1g};\quad Q_{xx}-Q_{yy},\ Q_{xy}\sim E_g;\quad Q_{xz},\ Q_{yz}\sim E_g.
\]
The component $Q_{zz}$ behaves as a scalar. Decomposing the direct products of
$E_g$ with all the group representations, we obtain the selection rules for
the nondiagonal matrix elements of the remaining components of $Q_{ik}$:
\[
 E_g\leftrightarrow A_{1g},A_{2g},E_g;\qquad E_u\leftrightarrow A_{1u},A_{2u},E_u.
\]
As follows from (4), diagonal elements can be nonzero only for states of type
$E_g$ and $E_u$.
